\section{Five terms at discriminant \texorpdfstring{$-131$}{-131}}\label{sec:quintic-sieve}

The class polynomial for $\Delta=-131$ has degree five. We show that
\emph{any} five distinct nonnegative powers of a root are linearly
independent over $\Q$. The test uses six primes. Each prime restricts the
possible exponent residues, including the case when two exponents have the
same residue.

\begin{lemma}[A finite test for unbounded exponents]\label{lem:confluent}
Let $H\in\Z[X]$ be monic and irreducible, let $x$ be the class of $X$
in $A=\Z_{(p)}[X]/(H)$, and suppose $x^m=a+pD$ in $A$ for a
rational integer $a$ prime to $p$. Fix distinct nonnegative integers
$e_1,\ldots,e_k$ and let $n_r$ count those congruent to $r$ modulo $m$.
If $p>k+1$ and the following columns are linearly independent in
$A/pA$ over $\mathbb F_p$, then $x^{e_1},\ldots,x^{e_k}$ are linearly
independent over $\Q$:
\[
                x^r D^j,\qquad 0\le r<m,\quad 0\le j<n_r .
\]
\end{lemma}
\begin{proof}
Within each residue block write $e=r+mt$ and divide the column by $a^t$.
It becomes $x^r(1+pD/a)^t$. For the first $j+1$ distinct values of $t$
in this block, form the Newton divided difference of order $j$ and
multiply by $j!/p^j$. Doing this successively gives an invertible
triangular change of columns over $\Q$.

Expand by the binomial theorem up to the largest $t$ in the block.
A divided difference of order $j$ annihilates $\binom{T}{i}$ for $i<j$
and equals $1/j!$ for $i=j$. For $i>j$ its denominator divides $i!$:
the polynomial $i!\binom{T}{i}$ has integer coefficients, and divided
differences of integer polynomials at integer arguments are integers.
The latter identity follows on monomials, whose order-$j$ divided
differences are the complete homogeneous symmetric polynomials of
degree $i-j$ at the chosen arguments.
For $i>j$ we have $i-j>v_p(i!)$. This is immediate for $i<p$.
For $i\ge p$ use $j\le k-1\le p-3$ and
$v_p(i!)\le (i-1)/(p-1)<i-j$.
Every transformed column therefore belongs to $A$, with reduction
$x^r(D/a)^j$ modulo $p$. Its scalar $a^{-j}$ is nonzero. Independence
after reduction implies independence over $\Q$ since $A$ is free over
$\Z_{(p)}$ on $1,x,\ldots,x^{\deg H-1}$. The triangular column operations
then give the conclusion.
\end{proof}

\begin{theorem}\label{thm:131}
The least support for discriminant $-131$ is six.
\end{theorem}
\begin{proof}
Use the certified class polynomial $H=H_{-131}$ in the accompanying data.
For each pair $(p,m)$ below, the checker reduces $x^m$ modulo
$(H,p^2)$. The constant coefficient is a $p$-adic unit, and every
other coefficient is divisible by $p$. Taking $a$ equal to the constant
coefficient determines $D\bmod p$.

Suppose five powers are dependent. Since $x\ne0$, divide by the lowest
power to obtain the exponent $0$ and four other exponents. At any
modulus $N$ we retain the sorted quadruple of the four residues, allowing
repetitions and zeros. On lifting $n\mid N$, add arbitrary multiples of
$n$ to the four old residues to obtain all representatives below $N$.
For equal old residues it suffices to take nondecreasing lifts among
those positions. Sorting the resulting quadruple gives every possible
lift exactly once. At each indicated prime, keep a quadruple only if the
five columns of Lemma~\ref{lem:confluent} are dependent. The distinguished
exponent $0$ supplies the first column $1$, so this test is a
four-by-four determinant after deleting its pivot row.
\begin{center}\small
\begin{tabular}{rrrr}
\toprule
$N$ & $(p,m)$ tested at this stage & Quadruples tested & Remaining\\
\midrule
200    & $(601,200)$                    &68,685,050 &119,264\\
1,400  & $(1051,350),\ (4201,1400)$      &279,218,674&196\\
4,200  & $(6301,1050)$                  &15,876     &5,876\\
12,600 & $(757,252)$                    &475,956    &214,120\\
25,200 & $(337,112)$                    &3,425,920  &0\\
\bottomrule
\end{tabular}
\end{center}
The modulus at every stage is divisible by each period used there.
Thus a genuine exponent set would give a surviving quadruple at every
stage. The last count proves independence for all five-element sets.
A relation on fewer powers would give a dependent five-element set by
adjoining distinct powers. Finally $H$ itself has at most six terms.
The polynomial reduction and determinant checks are reproduced by
\texttt{python3 checks/verify\_last\_quintic.py} from the archive root.
A separately written basis recurrence, cofactor computation and lift enumerator reproduce all five rows in \texttt{checks/independent\_quintic.py}.
\end{proof}

\paragraph{Exact minima through degree five.}
Combining the two finite censuses proves $s=h+1$ for every order of class number four or five. The same equality holds for $1\le h\le3$ except at $j=0$, where $s=1$. A nonzero rational $j$ has a binomial linear polynomial. For $h=2$ one has $D\ge15$, and the principal conjugate has larger absolute value than the other, since $e^{\pi\sqrt D}-2079>e^{\pi\sqrt D/2}+2079$. A binomial has all its nonzero complex roots on one circle. For $h=3$ the result is Theorem 6.6 of \cite{riffaut}, independent of GRH. Our checker \texttt{cm/verify\_all\_cubics.py} also certifies all 25 cubic orders using one disk per order, in characteristics $2,3,7,13,17$. It uses the degree-three instance of the conductor census above with $|d_K|\le907$, $f\le162$. Thus the complete table through $h=5$ is independent of hypotheses on zeros.
